%
% File: literature review
%
\chapter{Literature Review}
\label{chap:litreview}
\setlength{\parskip}{1em}
%
\section{Theoretical Foundations: CTML and CLT}

The Cognitive Theory of Multimedia Learning, developed by Richard Mayer over three decades, posits that meaningful learning occurs when learners actively integrate verbal and visual representations of the same content into coherent mental models. The theory rests on three assumptions inherited from cognitive psychology: the dual-channel assumption, that humans process auditory-verbal and visual-pictorial information through distinct channels; the limited-capacity assumption, that each channel can hold only a small number of elements at once; and the active-processing assumption, that learning requires deliberate cognitive effort to select, organise, and integrate information. From these assumptions Mayer derives a set of design principles --- the modality, redundancy, signalling, contiguity, coherence, and personalisation principles, among others --- that have been replicated across hundreds of experimental studies.

Cognitive Load Theory, articulated by John Sweller and colleagues, offers a complementary lens. CLT distinguishes intrinsic load (the inherent complexity of the material), extraneous load (the load imposed by the manner of presentation), and germane load (the load associated with the construction of schemas in long-term memory). Effective instruction minimises extraneous load and channels working-memory resources towards germane load. For learners with low language proficiency or with disabilities such as dyslexia, the linguistic surface of textual material constitutes a significant source of extraneous load: the cognitive cost of decoding language competes with the cognitive resources available for understanding the content itself.

The intersection of CTML and CLT yields a clear prediction: substituting or supplementing textual explanations with well-designed visual explanations should reduce extraneous load and free working memory for schema construction. \citet{lawson2024} provide recent experimental support for the active-processing component of this prediction, finding that summarising during multimedia animations improved learning outcomes while passive viewing did not, suggesting that the visual channel alone is insufficient unless paired with cognitive engagement.

\citet{twabu2025} propose an extension of these frameworks for the AI era. Their model introduces three additions: AI-mediated schema creation, in which a generative system scaffolds the learner's transition from novice to expert representations; adaptive cognitive load management, in which AI dynamically calibrates the complexity of the visual representation to the learner's measured load; and human--AI collaborative learning, in which the learner and the system iteratively refine a representation. The framework remains conceptual rather than empirically validated, but it offers a vocabulary for analysing AI-generated visual explanations as theoretically motivated rather than merely technologically convenient.

\section{AI Applications for Learners with Disabilities}

\citet{panjwani2024} provide the most comprehensive recent map of AI applications for students with learning disabilities, identifying seven categories across sixteen empirical studies: adaptive learning systems, facial expression recognition, chat robots, communication assistants, mastery learning environments, intelligent tutors, and interactive robots. Their analysis reveals two patterns of immediate relevance to the present review. First, adaptive learning systems dominate the empirical literature, while AI tools that produce or manipulate visual representations are relatively scarce. Second, the disability focus is heavily skewed: ten of the sixteen studies addressed dyslexia, with autism spectrum disorder a distant second and dyscalculia represented by only a single study. The authors classify the level of integration using the SAMR-LD model --- Substitution, Augmentation, Modification, Redefinition for learning disabilities --- and find that AI applications occupy all four levels, although the highest level (Redefinition, in which AI enables tasks otherwise impossible) remains rare.

\citet{drosos2025} extend this picture by reviewing 139 studies on AI, virtual reality, and large language models in special education. Their thematic analysis surfaces several application clusters of direct relevance: multimodal conversational systems for non-verbal communication, AI-generated images for augmentative and alternative communication, individualised educational support for children with developmental disorders, and VR-based social skills training. The authors document a clear historical shift from earlier assistive technologies, which substituted for impaired functions (text-to-speech, screen reading), towards generative AI tools, which produce new representational content tailored to the learner. This shift is consequential because it changes the question from ``how can the learner access the existing material?'' to ``what material best fits this learner?''

A particularly important strand of this literature concerns autism spectrum disorder. \citet{urrea2024} systematically review thirteen studies of technology-assisted vocabulary learning for children with ASD, including tablet applications, video modelling, computer-assisted instruction, and augmented reality. They report that visual systems such as the Picture Exchange Communication System remain a strong baseline, and that technology adds value primarily when it integrates structured visual scaffolding with explicit instruction. \citet{khoirunnisa2024} describe a complementary AR system for reading skills in ASD learners, in which visual overlays adapt to individual learner profiles. While neither study uses generative AI as such, both establish the population-level evidence base on which AI-generated visual explanations can build.

\citet{buhler2024} approach the same population from a different angle, conducting semi-structured interviews with students with disabilities in higher education about their use of generative AI. The participants --- many with ADHD, dyslexia, dyspraxia, or autism --- reported that ChatGPT and similar tools helped them structure ideas, overcome writer's block, and convert dense academic content into more accessible formats. Concerns centred on AI inaccuracy, academic integrity, and the cost of subscription tiers. The study is notable because it documents that learners themselves perceive generative AI as an assistive technology, validating the premise of the present review from the user side.

\section{Multimodal AI and Visual Generation in Education}

The technical capability that distinguishes the current moment from earlier waves of educational AI is multimodality. \citet{bewersdorff2024} develop the most explicit theoretical treatment of multimodal large language models in education, arguing that MLLMs change three things at once: they can interpret learner inputs in any modality, they can produce explanations in any modality, and they can move between modalities within a single interaction. The authors illustrate this with several scenarios drawn from science education: a diagram of a chemical reaction is augmented with a textual explanation calibrated to the student's reading level; a student's hand-drawn sketch of cellular respiration is corrected and re-rendered as a clean diagram; a multimodal assessment combines a generated image with a written prompt. Crucially, \citet{bewersdorff2024} ground each scenario in CTML, showing that the design choices follow directly from established multimedia principles rather than from technological novelty.

\citet{osman2024} operationalise a similar approach in the form of an AI educational video assistant for higher education. Their tool comprises three modules: a Transcription Module that converts lecture audio to text using OpenAI Whisper, an Engagement Module that uses a large language model to support interactive question-and-answer, and a Reinforcement Module for spaced review. Each module is mapped explicitly to specific CTML principles. The Transcription Module supports the modality principle by allowing learners to switch between auditory and textual representations. The Engagement Module supports the personalisation principle by responding in conversational language. The Reinforcement Module supports the spaced-practice principle. Both human evaluators and automated metrics confirmed gains in learning outcomes, although the study did not specifically isolate effects for learners with disabilities or low language proficiency.

A complementary line of work examines the conditions under which visual content actually produces learning. \citet{lawson2024} conducted an experimental study in which students viewed annotated animations about greenhouse gases and were asked to summarise, draw, or simply watch. Summarising produced significant learning gains; drawing did not. The result is a useful corrective to a naive enthusiasm for visual content: a generated diagram, however accurate, does not produce learning if the learner is a passive consumer of it. The pedagogical value of AI-generated visual explanations therefore depends on the activity structure that surrounds them.

\citet{ahmad2025} bring these strands together in a review that explicitly addresses the dual challenge of disability support and language barriers. They document AI-driven adaptive learning systems, intelligent tutoring platforms such as ALEKS and Q-interactive, NLP-based text simplification, bilingual glossary generation, and assistive technologies including text-to-speech, closed captioning, and image recognition. Their headline empirical finding --- that AI-enabled hybrid learning models have produced improvements of up to twenty-five per cent in vocabulary learning and thirty per cent in reading comprehension in certain contexts --- provides the strongest available evidence for the population of interest, although the authors caution that effects are highly context-dependent and that ethical concerns about bias and privacy remain unresolved.

\section{Identified Gap}

Three observations emerge from this body of work. First, the theoretical apparatus exists: CTML and CLT, with the recent extensions proposed by \citet{twabu2025}, provide a coherent vocabulary for analysing AI-generated visual explanations. Second, the empirical literature on AI in special education and on multimodal AI in mainstream education is growing rapidly, but the two literatures rarely meet. Studies of AI for learners with disabilities tend to describe assistive applications without engaging multimedia learning theory; studies of multimodal AI in education tend to describe theoretically grounded applications without examining differential effects for learners with low language proficiency or disabilities. Third, where the two literatures do meet --- for example in \citet{ahmad2025} and \citet{drosos2025} --- the treatment is broad rather than focused, surveying many tools without close examination of how visual generation specifically supports the populations of interest.

The gap is therefore not an absence of research but a lack of integration. A systematic review oriented specifically to AI-generated visual explanations, framed by CTML and CLT, and focused on learners with low language proficiency or learning disabilities, would synthesise what is currently known, expose the conditions under which the technology supports learning, and identify which questions remain open. The methodology adopted to address this gap is set out in the next chapter.
